\appendix
\section{Vanilla Partitioning Problem}

\subsection{Problem Definition}

We are given a graph $G=\left(V,E\right)$ with $n$ nodes and $m$ edges. Each node has two node weights $h_i$ and $s_i$ to denote hardware and software costs.
The edges denote communication between the nodes and have associated communication cost $c_{ij}$ for the edge $e$ connecting nodes $i$ and $j$.
The aim is to partition the nodes $V$ into non-overlapping sets $V_S$ and $V_H$ denoting software and hardware nodes respectively; $V=V_H \cup V_S$ and $V_H \cap V_S = \emptyset$.
We denote the set of edges between the two sets by $E_{\cP}$. 
Note that we will only consider communication cost for these edges which are between hardware and software nodes.
The hardware cost for a partition $\cP$ is defined as sum of hardware cost for the nodes belonging to the hardware nodes 
$H_{\cP} = \sum_{v\in V_H}h_v$. Similarly, the software cost for the partition is $S_{\cP} = \sum_{v\in V_S}s_v$.
The communication cost is defined as $C_{\cP} = \sum_{e:(i,j) \in E_\cP} c_{ij}$.

\subsection{Optimization Formulation}
We define binary variable $x_i$ for each node $i \in V$; $x_i=1$ if node $i$ is a software node in partitioning $\cP$, else $x_i=0$ for a hardware node.
Additionally, we denote the edge weighted adjacency matrix for the graph $G$ to be $\bC$, where the element along the $i^{th}$ row and $j^{th}$ column $C_{ij}$ is \[C_{ij}=\begin{cases}c_{ij},~ \textrm{if } e:(i,j) \in E \\ 0,~ \textrm{otherwise}\end{cases}\] Note that $\bC$ is symmetric.
The communication cost for the partition $\cP$ is 
\[C_{\cP} = \sum_{e:(i,j) \in E_\cP} c_{ij}=\sum_{i\in V}\sum_{j\in V}x_i C_{ij} (1-x_j)\]
We can stack the entries $h_i$, $s_i$ and $x_i$ in $n$-length vectors $\bh$, $\bs$ and $\bx$ respectively. For the partition $\cP$, the hardware cost is $H_{\cP}=\bh^T\left(\ones-\bx\right)$, the software cost is $S_{\cP}=\bs^T\bx$, and the communication cost $C_\cP=\left(\ones-\bx\right)\bC\bx$.
The hardware-software partitioning problem can be stated as
\begin{align*}
    \cP^\star = \arg \min & ~ H_\cP\\
    \textrm{s.t.} & ~ S_\cP + C_\cP \leq R_0
\end{align*}
which we can write as
\begin{subequations}
\begin{align}
    \bx^\star = \arg \min & ~ \bh^T\left(\ones - \bx\right)\\
\textrm{s.t.} & ~ \bs^T\bx + \left(\ones - \bx\right)^T\bC\bx \leq R_0\\
& ~ x_i \in \{0,1\},~\forall i=1,2,\cdots,n
\end{align}
\end{subequations}
\begin{relaxation}
    We relax the integer constraint $x_i \in \{0,1\}$ to $0 \leq x_i \leq 1,~ \forall i=1,2,\cdots,n$. 
\end{relaxation}
The relaxed optimization problem is a quadratic constrained linear problem
\begin{subequations}
\begin{align}
    \bx^\star = \arg \min & ~ \bh^T\left(\ones - \bx\right)\\
\textrm{s.t.} & ~ \left(\bs+\bC\ones\right)^T\bx - \bx^T\bC\bx \leq R_0\\
& ~ \zeros \leq \bx \leq \ones
\end{align}
\end{subequations}
Note that $\bC \succeq 0$, and hence, the constraint $\left(\bs+\bC\ones\right)^T\bx - \bx^T\bC\bx \leq R_0$ is non-convex.
We now define $(n+1)\times(n+1)$ matrices
\[\bX=\begin{bmatrix}\bx\bx^T & \bx \\ \bx^T & 1\end{bmatrix} \textrm { and } \bH=\begin{bmatrix}0 & \frac{\bh}{2} \\ \frac{\bh^T}{2} & 0\end{bmatrix}\]
to pose the above problem as a semidefinite program (SDP). Using the linearity and cyclic properties of \emph{trace} of a matrix, we have
\begin{subequations}
    \begin{align}
        & \tr{\bH\bX} = \bh^T\bx\\
        & \tr{\begin{bmatrix}\bC&-\frac{\bs+\bC\ones}{2}\\-\frac{\left(\bs+\bC\ones\right)^T}{2}&R_0\end{bmatrix}\bX} = \bx^T\bC\bx - \left(\bs + \bC\ones\right)^T\bx + R_0
    \end{align}
\end{subequations}
Additionally, we want to add some more constraints to make the problem more tight. These include
\begin{subequations}
    \begin{align}
        \forall i=1,2,\cdots,n \quad &0 \leq x_i \leq 1 ~~ \Rightarrow 0 \leq \be_i^T\bx \leq 1\\
        \forall i=1,2,\cdots,n \quad &X_{ii} = x_i ~~ \Rightarrow \bx^T\bE_{ii}\bx = \be_i^T\bx\\
        \forall i,j=1,2,\cdots,n \quad &X_{ij} \leq x_i ~~ \Rightarrow \bx^T\bE_{ij}\bx \leq \be_i^T\bx\\
        \forall i,j=1,2,\cdots,n \quad &X_{ij} \leq x_j ~~ \Rightarrow \bx^T\bE_{ij}\bx \leq \be_j^T\bx\\
        \forall i,j=1,2,\cdots,n \quad &X_{ij} \geq 0 ~~ \Rightarrow \bx^T\bE_{ij}\bx \geq 0
    \end{align}
    \label{eq:tight}
\end{subequations}
Here, $\bE_{ij}$ denotes a $n\times n$ matrix with all zeros, except the element along the $i^{th}$ row and $j^{th}$ column being $1$, and the vector $\be_i$ denotes the $i^{th}$ column of the $n\times n$ identity matrix.
Using the linearity and cyclic properties of \emph{trace} of a matrix, we have the following identities
\begin{subequations}
    \begin{align}
        &  \tr{\begin{bmatrix}\zeros&\frac{\be_i}{2}\\\frac{\be_i^T}{2}&0\end{bmatrix}\bX} = \be_i^T\bx\\
        & \tr{\begin{bmatrix}\bE_{ij}&-\frac{\be_i}{2}\\-\frac{\be_i^T}{2}&0\end{bmatrix} \bX} = \bx^T\bE_{ij}\bx - \be_i^T\bx
    \end{align}
    \label{eq:identities}
\end{subequations}
The SDP problem is
\begin{subequations}
\begin{align}
\textrm{minimize} & ~ \bh^T\ones - \textrm{Tr}(\bH\bX)\\
\textrm{subject to} & ~ \tr{\begin{bmatrix}\bC&-\frac{\bs+\bC\ones}{2}\\-\frac{\left(\bs+\bC\ones\right)^T}{2}&R_0\end{bmatrix}\bX}\geq 0\\
& ~ 0\leq \tr{\begin{bmatrix}\zeros&\frac{\be_i}{2}\\\frac{\be_i^T}{2}&0\end{bmatrix}\bX}\leq 1,~~ \tr{\begin{bmatrix}\bE_{ii}&-\frac{\be_i}{2}\\-\frac{\be_i^T}{2}&0\end{bmatrix}\bX}= 0, ~~ \forall i\\
& ~ \tr{\begin{bmatrix}\bE_{ij}&-\frac{\be_i}{2}\\-\frac{\be_i^T}{2}&0\end{bmatrix}\bX}\leq 0,~~\tr{\begin{bmatrix}\bE_{ij}&-\frac{\be_j}{2}\\-\frac{\be_j^T}{2}&0\end{bmatrix}\bX}\leq 0 ~~\forall i,j\\
& ~ \tr{\begin{bmatrix}\bE_{ij}&\zeros\\\zeros^T&0\end{bmatrix}\bX}\geq 0~~\forall i,j, ~~ \bX \succeq 0\\
& ~ \textrm{rank}(\bX) = 1
\end{align}
\end{subequations}

\section{Scheduling Constrained Partitioning Problem}

\subsection{Problem Definition}
We are given a graph $G=\left(V,E\right)$ with $n$ nodes and $m$ edges. 
Each node has weights $h_i$ and $s_i$ to denote hardware and software execution times.
Also, each node has an associated hardware area cost $a_i$.
The edges denote communication between the nodes and have associated communication cost $c_{ij}$ for the edge $e$ connecting nodes $i$ and $j$.
The communication delay is taken into account only when the adjacent nodes belong to different partitions.
The aim is to partition the nodes $V$ into non-overlapping sets $V_S$ and $V_H$ denoting software and hardware nodes respectively; $V=V_H \cup V_S$ and $V_H \cap V_S = \emptyset$.
We denote the set of edges between the two sets by $E_{\cP}$. 
Additionally, we need to ensure that software nodes are executed sequentially, while hardware nodes can execute in parallel if their dependent nodes are executed.
Our objective is to minimize the overall execution time given a limited available area for hardware nodes.

\subsection{Optimization Formulation}

\noindent\textbf{Decision variables.}~We define binary variable $x_i=0$ if node $i$ is assigned to hardware partition, $x_i=1$ if assigned to software partition.
We also define binary variable $z_{ij}$ for each edge $e:=\left(i,j\right)\in E$connecting nodes $i$ and $j$. We denote $z_{ij}=1$ if the nodes $i$ and $j$ belong to different partitions.
Let $t_i$ and $f_i$ denote the start and end time of the node $i$. The total execution time of the task graph is denoted by $T$.
We also define binary variable $y_{ij}$ for each pair of nodes $i,j$, $i \neq j$.
This binary variable is used to ensure that if the pair of nodes $i,j$ are assigned to software partition, they need to be executed sequentially; i.e., if $y_{ij}=1$, then node $i$ precedes node $j$, and for $y_{ij}=0$, the opposite happens.
However, we need to ensure that this sequential constraint is activated for only a pair of software nodes.

\noindent\textbf{Hardware area constraint.}~The overall hardware area is constrained by the maximum allowable hardware area $A_{\textrm{max}}$.
\begin{equation}
    \sum_{i=1}^{n}a_i\left(1-x_i\right) \leq A_{\textrm{max}}
    \label{eq:hardware-area}
\end{equation}

\noindent\textbf{Binary edge variable.}~The binary variable $z_{ij}$ and the partition variable $x_i,x_j$ corresponding to the adjacent nodes are related by
\begin{subequations}
    \begin{align}
        &z_{ij} \geq x_i - x_j, \quad &\forall \left(i,j\right) \in E\\
        &z_{ij} \geq x_j - x_i, \quad &\forall \left(i,j\right) \in E\\
        &z_{ij} \leq x_i + x_j, \quad &\forall \left(i,j\right) \in E\\
        &z_{ij} \leq 2 - x_i - x_j, \quad &\forall \left(i,j\right) \in E
    \end{align}
    \label{eq:edge-selection}
\end{subequations}

\noindent\textbf{Execution time constraints.}~The execution time of each node is denoted 
\begin{subequations}
    \begin{align}
        & T \geq f_i, \quad &\forall i \in V\\
        & f_i = t_i + h_i\left(1-x_i\right) + s_ix_i, \quad &\forall i\in V\\
        & f_i + c_{ij}z_{ij} \leq t_j, \quad &\forall \left(i,j\right) \in E
    \end{align}
    \label{eq:exec-time}
\end{subequations}

\noindent\textbf{Sequential execution of software nodes.}We define the following pair of constraints with $M$ being a large positive number.
\begin{subequations}
    \begin{align}
        & f_i \leq t_j + M \left(1 - y_{ij}\right) + M \left(2 - x_i -x_j\right), \quad &\forall i,j \in V, i \neq j\\
        & f_j \leq t_i + M y_{ij} + M \left(2 - x_i -x_j\right), \quad &\forall i,j \in V, i \neq j
    \end{align}
    \label{eq:sw-seq}
\end{subequations}
Note that, the constraints are activated only when $x_i=x_j=1$ or the nodes $i,j$ are both software nodes, when the term $2-x_i-x_j=0$.
For other cases, the constraints are not active since $M$ is sufficiently large.
Once the constraints are activated, the value of $y_{ij}$ denotes the precedence of the software node operations.

The overall optimization problem is
\begin{subequations}
    \begin{align}
        \textrm{minimize} & ~ T\\
        \textrm{subject to} & ~ \left(\ref{eq:hardware-area}\right),\left(\ref{eq:edge-selection}\right),\left(\ref{eq:exec-time}\right),\left(\ref{eq:sw-seq}\right)
    \end{align}
\end{subequations}